\section{Building and finding: the scripts}\label{sec:build-and-find}

One executable builds; a small family searches. \file{build-pdfs} compiles
course documents (with their views, and with output routed where you want
it). The searchers are per-kind faces of one shared engine:
\file{find-problems} searches the problem bank, \file{find-facts} the fact bank
and \file{find-passages} the passage bank -- each a two-line shim that hands the
right \texttt{-{}-kind} to \file{find-meta}, where all the real logic lives once
-- and their interactive twins \file{pick-problems} / \file{pick-facts} /
\file{pick-passages} (shims over \file{pick-meta} the same way) thread those
searches through fzf. Every one of them prints its full reference with
\texttt{-{}-help}.

\subsection{\texorpdfstring{\texttt{build-pdfs}}{build-pdfs}}

\begin{manexsource}
$ build-pdfs              # compile every .tex in this directory
$ build-pdfs test.tex     # just this one (FILE.tex -> FILE.pdf)
\end{manexsource}

The script runs latexmk from each file's own directory, so relative imports
resolve and the directory's \file{.latexmkrc} puts the machinery on
\env{TEXINPUTS}. It installs that \file{.latexmkrc} on the way in: latexmk
reads only the rc in its startup directory and never walks up, so before
compiling, \file{build-pdfs} copies \file{latexmkrc-shared} from the machinery
into the source's directory whenever the file there is missing or out of date.
A new topic folder therefore needs no setup, and an edit to the master reaches
every course on its next build. A copy whose \texttt{managed by build-pdfs}
marker line has been deleted counts as hand-tuned and is never overwritten.
Auxiliary files (\file{.aux}, \file{.log}, \dots) go to a
per-directory folder under \file{~/.cache/latexmk/} -- they clutter neither
the repo nor the output directory -- keyed by the directory's repo-relative
path, so two topics' same-named \file{problems/} folders never share
latexmk state.

\subsubsection{Views: one source, several PDFs}

A magic comment in a source's \emph{leading} comment block (it must precede
\cs{documentclass}) requests extra PDFs:

\begin{manexsource}
%! views: solutions, hints, instructor
\documentclass[12pt,thursday]{Assessment}
\end{manexsource}

Four views are wired, each injecting a layer definition ahead of the source
so the call site needs no class option:

\begin{itemize}
  \item \env{solutions} $\to$ \file{FILE-SOLUTIONS.pdf}: defines
    \cs{ShowSolutions}, so \env{solution} bodies and the boxed MC keys
    print.
  \item \env{hints} $\to$ \file{FILE-HINTS.pdf}: defines \cs{ShowHints}.
    Independent of solutions -- a hint is student-facing help, so list both
    views if you want both PDFs.
  \item \env{projector} $\to$ \file{FILE-PROJECTOR.pdf}: defines
    \cs{ProjectorView}, switching \file{LectureNotes} to the 17\,pt room
    copy.
  \item \env{instructor} $\to$ \file{FILE-INSTRUCTOR.pdf}: the staff
    bundle -- solutions $+$ hints $+$ instructor notes, inside the
    page-edge frame. A view is a \emph{bundle of layers}; the instructor
    copy is the staff copy (TAs get it too -- there is no separate marker
    tier).
\end{itemize}

An unrecognised view name warns and is skipped; the student PDF always
builds. Under the default overlay mode, every view of a document has
identical page breaks (section~\ref{sec:overlay}).

\subsubsection{Where the PDFs go}

The output mirror is one \file{.MIRRORDIR} file, found by walking up from the
build directory to the nearest ancestor that holds one (the repo top level,
the way \file{.latexmkrc} finds the machinery). Its first bare line names a
base directory (leading \texttt{\~{}} expanded) that \emph{mirrors the repo
tree}: a source built in \file{REPO/integrals/lectures} ships its PDF to
\file{BASE/integrals/lectures}. The mirror subdirectory is created if absent,
but the base itself must exist -- if the cloud drive is unmounted, the script
warns and builds locally rather than scattering a stray tree. By default the
PDF goes \emph{straight there}, with no local copy; add a line reading
\texttt{BOTH} to build locally \emph{and} copy. Blank and \texttt{\#}-comment
lines are ignored; the base path and the \texttt{BOTH} keyword may come in
either order. A nearer \file{.MIRRORDIR} in a subdirectory wins, letting a
subtree mirror somewhere else. With no \file{.MIRRORDIR} above it at all, the PDF
lands next to its source.

\file{SAMPLE-MIRRORDIR}, shipped with the machinery, is the commented template --
copy it to \file{.MIRRORDIR} and uncomment. \file{.MIRRORDIR} is meant to be
git-ignored: each machine points its own way.

\paragraph{Routing a view to a named folder.}
Sometimes one view of a document belongs somewhere else too -- the instructor
copy of a test in a folder shared with TAs, say, while the rest of the build
still mirrors to your private cloud. Add \texttt{LABEL: path} lines to
\file{.MIRRORDIR} and tag the view with a matching \texttt{:LABEL} suffix in the
magic comment:

\begin{manexsource}
%! views: solutions, instructor:TA-dir
\end{manexsource}

\noindent That drops \file{FILE-INSTRUCTOR.pdf}, flat, into the directory
\texttt{TA-dir:} names \emph{in addition} to the mirror; \env{solutions} still
goes only to the mirror. The labelled directory must already exist (a missing
label or directory warns and skips just that extra copy -- the build still
succeeds and still mirrors). Because the source names only a \emph{label} and
the git-ignored \file{.MIRRORDIR} maps it to a path, no machine-specific path ever
enters a tracked source. The pseudo-view \env{default} names the base (student)
PDF itself, so \texttt{default:LABEL} routes it the same way; a bare
\env{default} is a harmless no-op.

Routing decides \emph{where} a view goes, not \emph{when} it is released: if
you would rather TAs not see solutions until after the assessment, point the
label at a private staging folder and copy across by hand at release time --
\file{build-pdfs} deliberately knows nothing about dates.

\subsection{\texorpdfstring{\texttt{find-problems} / \texttt{find-facts} /
  \texttt{find-passages}}{find-problems / find-facts / find-passages}}

\begin{manexsource}
$ find-problems [--topic PAT] [--type PAT] [--difficulty PAT]
                [--source PAT] [--untagged] [--vocab] [--import]
                [--tsv] [PATH ...]
$ find-facts    [--topic PAT] [--source PAT] [--untagged] [--vocab]
                [--import] [--tsv] [PATH ...]
$ find-passages [--audience PAT] [--untagged] [--vocab]
                [--import] [--tsv] [PATH ...]
\end{manexsource}

Reads the \cs{ProblemMeta} (resp.\ \cs{FactMeta}, \cs{PassageMeta}) block off
every \file{.tex} file under each \env{PATH} (default: the current directory,
recursively) and reports the matches -- without running LaTeX, which is the
point (section~\ref{sec:problem-files}). With no filters, it lists every
file that has a block. The kind fixes the block macro it reads, the
\cs{importproblem}~/ \cs{importfact}~/ \cs{importpassage} lines it emits, and
the filter set -- facts carry no \env{type} or \env{difficulty}, and a passage is
searched by \env{audience} alone, so the three have four, two and one filters
respectively. Everything below holds for all three.

One difference is worth stating here rather than leaving to be discovered:
\file{find-passages} accepts no \texttt{-{}-with}. A passage has no injectable
keys, because the only option \cs{importpassage} takes is \env{level}, and where
a passage's natural depth lands belongs to the importing document rather than to
the passage (section~\ref{sec:ref-passagemeta}). Emitted passage lines are
therefore always bare. The passage bank also sits in a sibling submodule rather
than a directory of the course, so give it as an operand ---
\texttt{find-passages teaching/passages} --- or a bare run from the course root
will walk everything else too.

Filters combine with AND. \env{PAT} is an extended regular
expression matched case-insensitively against each comma-separated item, so
\texttt{-{}-topic power} finds \texttt{power-rule} and
\texttt{-{}-topic '\^{}power-rule\$'} matches exactly. A filter on a key a
file does not set never matches -- \texttt{-{}-difficulty basic} will not
sweep in untagged files.

Four modes round it out:

\begin{itemize}
  \item \texttt{-{}-vocab} ignores the filters and prints the type and topic
    vocabulary in use, with counts -- the antidote to a bank where half the
    problems say \emph{integration} and half say \emph{integrals}. Read it
    before inventing a type.
  \item \texttt{-{}-untagged} inverts: files with \emph{no} metadata block
    at all -- the tagging backlog.
  \item \texttt{-{}-import} prints the matches as paste-ready
    \cs{importproblem} (or \cs{importfact}) lines. Two opt-in modifiers
    refine the output:
    \texttt{-{}-with LIST} injects a comma subset of the injectable keys as
    options where a file carries them -- \env{number} and \env{name} for a
    problem; for a fact also \env{kind} and \env{name}, which choose and
    title the box it is wrapped in (\env{points} stays manual --
    it is extrinsic) -- and \texttt{-{}-from DIR} rewrites each line's directory
    argument relative to \env{DIR}, the importing document's folder, so the
    pasted line needs no \texttt{../} hand-fixing.
  \item \texttt{-{}-tsv} prints one tab-separated row per match (path, then the
    metadata) -- the machine-readable feed behind the picker below. An operand
    may now be a single \file{.tex} file, not just a directory, so a picker can
    ask \texttt{-{}-import} about exactly the files it selected.
\end{itemize}

The parser is textual and deliberately simple, which is where the
one-key-per-line discipline of section~\ref{sec:problem-files} comes from:
a reflowed block still compiles, but this script no longer sees it.

\subsection{Picking from a bank interactively}

\file{pick-problems} (and \file{pick-facts} / \file{pick-passages}, identically)
threads those two
machine modes through
\href{https://github.com/junegunn/fzf}{fzf}: it feeds \texttt{-{}-tsv} into the
picker (fuzzy-search the metadata, \file{bat}-preview the \file{.tex},
multi-select with Tab), then pipes the chosen files back through
\texttt{-{}-import} so the accepted set comes out as \cs{importproblem} (or
\cs{importfact}, \cs{importpassage}) lines on
standard output. Its \texttt{-{}-from} and any \file{find-problems} filters pass
straight through, e.g.\ \texttt{pick-problems -{}-from \_assessment -{}-type
multiple-choice derivatives/}. The Vim commands \texttt{:ImportProblem}
(\texttt{<Space>ip}), \texttt{:ImportFact} (\texttt{<Space>if}) and
\texttt{:ImportPassage} (\texttt{<Space>ig}) are the same
pipeline, inserting the lines at the cursor
with \texttt{-{}-from} pinned to the buffer's own directory -- so picking a
problem while writing an assessment drops in a ready line, needing only its
\env{points}, and picking a passage drops in one needing nothing at all.

Two things about picking \emph{passages} in particular. The emitted lines are
separated by a blank line, matching how every call site writes them, which makes
each one a paragraph object --- and that matters because fzf returns a
multi-selection in \emph{input} order, not the order you ticked them, while a
passage's position in a document can be load-bearing wherever one passage refers
to another. So reorder after pasting: \texttt{dap}, \texttt{p}. Second, \env{level}\texttt{=subsection} is
always added by hand and is always safe to try: \file{CourseDocument} refuses a
shift that would push a heading past \cs{subsubsection}, loudly, at build time.
